% \vspace{-5pt}
\section{Introduction}
\label{dyn_sec:intro}

Reconstructing dynamic targets from images is a long-standing computer vision problem, with applications in motion analysis~\cite{liu2023building,chao2025part}, AR/VR~\cite{yan2025instant}, and dynamic scene understanding~\cite{wang2023root}. While recent progress in neural~\cite{pumarola2021d,cao2023hexplane,uzolas2023template} and Gaussian-based representations~\cite{Wu_2024_CVPR,wan2024template,huang2024sc,yao2025riggs} have shown impressive results, most methods rely on monocular or multi-view videos with dense temporal coverage, where rich motion cues and correspondences are available. 
% \vtxt{what does ``densely captured monocular video" mean?}

In real-world scenarios, however, such dense observations are not always accessible. For example, surveillance cameras often capture moving objects sparsely over time, especially in cluttered environments.
Moreover, when multiple cameras are available, their viewpoints can differ drastically, and the observed targets may exhibit significant motion and self-occlusion between observations.
Under this setting, temporal correspondences are difficult to establish, as appearance can change dramatically across sparse observations, making dynamic reconstruction highly ill-posed.

In this paper, we address this challenging setting of articulated dynamic reconstruction from sparse temporal observations, where only a few posed images from arbitrary viewpoints are available as illustrated in \figref{dyn_img:teaser}.
To solve this highly ill-posed problem, we consider a setting where we have access to additional structural information. 
% Initially, we start with the assumption that a rough skeleton graph and an initial static reconstruction at the first frame are available. 
Initially, we assume that a rough skeleton graph and a static reconstruction at the first frame are available. 
The initial reconstruction can be can be obtained from a standard multi-view setup ~\cite{schoenberger2016sfm, schoenberger2016mvs, kerbl3Dgaussians}.
Later on in \secref{dyn_sec:diffusion}, we will show how this assumption can be relaxed with a pre-trained generative model~\cite{liu2023zero, shi2023mvdream, tang2023dreamgaussian} using only a single image. 
Despite this additional information, the task remains difficult as the inputs do not yield a complete rigged model---the skeleton annotation can be noisy and contains only node positions and connectivity, while the joint poses, skinning weights, and point-to-part associations remain unknown. 

We present \ourmethod{} which, given the input skeleton graph and initial static reconstruction, learns a skeleton-driven deformation field that models coherent motion under sparse supervision. Our deformation field consists of a coarse skeleton joint pose estimator and a module that models fine-grained motion deformations. By allowing only the joint pose estimator to be time-dependent, our model enables smooth test-time motion interpolation while preserving learned local deformation details.
Experiments demonstrate that state-of-the-art (SOTA) dynamic reconstruction methods degrade significantly in this sparse setting, while \ourmethod{} achieves better reconstruction quality. 
% Specifically, our approach outperforms the baselines by up to 34\% in PSNR on synthetic datasets with sparse observations, and achieves comparable performance to dense monocular video methods on real-world datasets.
Furthermore, we show that the need for multi-view initialization can be relaxed using a diffusion-based generative prior, enabling dynamic reconstruction in real-world scenarios. 
Our contributions can be summarized as follows. 
\begin{compactitem} 
\item  We perform articulated dynamic reconstruction from sparse temporal observations, where only a few frames from arbitrary viewpoints are available.
\item We present a skeleton-driven deformation field that enables smooth motion interpolation under sparse supervision, and demonstrate that a pre-trained diffusion prior can be incorporated to fill in missing information.
\item Experiments show that our method outperforms SOTA methods by up to 34\% in PSNR on synthetic datasets with sparse observations, and achieves comparable performance to dense monocular video methods on real-world datasets with significantly fewer frames.
\end{compactitem}



\begin{comment}
\par
$\bullet$ Motivate the problem of dynamic target reconstruction from image observations

\par
$\bullet$ Most existing methods assume densely captured monocular videos with dense temporal motion cues and correspondences. 

\par
$\bullet$ However, in real-world scenarios, such dense observations might not always be available. For example, surveillance camera capturing a dynamic target in a cluttered scene. In this case, the observations are usually sparse across the temporal dimension, also the viewpoints can vary a lot between consecutive observations. Therefore, building dense correspondences for dynamic objects in this setup can be challenging. 

\par
$\bullet$ We aim for this challenging setting where we are interested in reconstructing the dynamic target from sparse temporal observations.


\par
$\bullet$ To solve this highly ill-posed problem, we assume access to a rough skeleton graph and an initial static reconstruction at the first frame, which can be obtained from a static reconstruction methods or potentially from a pre-trained generative model. 

\par
$\bullet$ Emphasize that this is still a difficult problem even when given skeleton graph and initial static reconstruction. The input does not yield a complete rigged model as the skeleton graph annotation can be noisy and contains only node's 3D positions and connectivity. The actual joint poses, skinning weights, and point-to-part association is still unknown. 
Moreover, with sparse observations, building dense correspondences is challenging as the viewpoint and object poses are different between observations. 
The goal is to infer how the target deforms and articulates over time under sparse observations.

\par
$\bullet$ We show in the experiment section that existing methods suffer when only sparse observation is available even with multi-view ground truth at the initial static state. We further demonstrate the potential of replacing the multi-view initialzation with a pre-trained generative model. 

\end{comment}